%% kapa: exam=1 round=29 year=2018 period=2 form=A subject=appraisal_law subject_ko=감정평가관계법규 seq=38 number=38
%% source: Q-Net 감정평가사 기출문제 (cst00309) artlSeq=5202808 file=extracted/1008/5202808/2018년 제29회 감정평가사 1차시험 기출문제/2018년 제29회 감정평가사 1차시험 기출문제/2교시 A형.pdf
%% conversion: kordoc 4.24.51 -> tools/gen1.py (auto); review=proofread-sonnet; flags=none
\begin{question}{38}
도시 및 주거환경정비법령상 가로주택정비사업에 관한 설명으로 옳지 않은 것은?
\begin{choices}
  \item 가로주택정비사업은 사업시행자가 가로구역에서 인가받은 관리처분계획에 따라 주택 등을 건설하여 공급하거나 보전 또는 개량하는 방법으로 한다.
  \item 가로주택정비사업의 사업시행자는 가로구역에 있는 기존 단독주택의 호수와 공동주택의 세대 수를 합한 수 이상의 주택을 공급하여야 한다.
  \item 특별자치도지사는 가로주택정비사업의 시행을 위하여 도시ㆍ주거환경정비기본계획을 수립하여야 한다.
  \item 가로주택정비사업은 조합이 시행하거나 조합이 조합원 과반수의 동의를 받아 시장ㆍ군수와 공동으로 시행할 수 있다.
  \item 가로주택정비사업의 경우에는 조합설립추진위원회를 구성하지 아니한다.
\end{choices}
\end{question}
\answer{3}
